% ============================================================================
% PHAN 3: CHUONG 2, CHUONG 3, KET LUAN VA TAI LIEU THAM KHAO
% Du an: Thiet ke bo ma hoa - giai ma Reed-Solomon RS(15,11) tren DE10-Lite
% ============================================================================

\chapter{Giới thiệu về Board FPGA Terasic DE10-Lite}
\label{ch:de10-lite}

Để hiện thực hóa và kiểm chứng thực nghiệm bộ mã hóa -- giải mã Reed--Solomon $\mathrm{RS}(15,11)$ trên phần cứng khả trình, đồ án lựa chọn kit phát triển \textbf{Terasic DE10-Lite} trang bị vi mạch FPGA \textbf{Intel MAX 10} mã \texttt{10M50DAF484C7G}. Chương này trình bày kiến trúc tổng quan của bo mạch, đặc tính tài nguyên của vi mạch MAX 10, nguyên lý điện -- logic của các ngoại vi giao tiếp và bảng quy hoạch gán chân kèm ràng buộc thời gian tĩnh cho toàn hệ thống.

\section{Tổng quan kiến trúc Kit Terasic DE10-Lite}
\label{sec:de10-overview}

Kit \textbf{Terasic DE10-Lite} là nền tảng FPGA phổ thông hướng tới thiết kế logic số và xử lý tín hiệu. Bo mạch tích hợp sẵn mạch nạp \textbf{USB-Blaster II} qua cổng USB Type-B, vừa cấp nguồn 5\,V cho toàn hệ thống (qua các bộ ổn áp hạ xuống $3{,}3\text{ V}$, $2{,}5\text{ V}$ và nguồn lõi $1{,}20\text{ V}$), vừa truyền tệp cấu hình \texttt{ReedSolomon.sof} trực tiếp vào FPGA qua giao thức JTAG mà không cần mạch nạp rời. Hệ thống tạo xung nhịp sử dụng bộ dao động thạch anh 50\,MHz cấp cho chân \texttt{MAX10\_CLK1\_50} (\texttt{PIN\_P11}) làm miền xung nhịp đơn duy nhất (\textit{single clock domain}). Về ngoại vi giao tiếp người dùng, bo mạch cung cấp 10 switch trượt \texttt{SW[9:0]}, 2 nút nhấn có mạch \textit{Schmitt Trigger} \texttt{KEY[1:0]}, 10 LED đơn \texttt{LEDR[9:0]} và 6 LED 7 đoạn \texttt{HEX5..HEX0} (cùng header mở rộng \texttt{JP1}, SDRAM và VGA chưa sử dụng trong phiên bản hiện tại). Hình~\ref{fig:de10-overview} mô tả sơ đồ khối kết nối giữa vi mạch MAX 10 và các ngoại vi phục vụ hệ thống $\mathrm{RS}(15,11)$.

\begin{figure}[!htb]
    \centering
    \resizebox{0.96\textwidth}{!}{%
    \begin{tikzpicture}[>=stealth, font=\footnotesize]
        % Khung FPGA trung tam
        \draw[thick, rounded corners=4pt, fill=blue!4] (4.8, 0.2) rectangle (10.3, 5.9);
        \node[font=\small\bfseries, align=center] at (7.55, 5.45) {Intel MAX 10 FPGA\\(\texttt{10M50DAF484C7G})};

        % Khoi thiet ke ben trong FPGA
        \draw[thick, rounded corners=3pt, fill=white] (5.25, 1.4) rectangle (9.85, 4.7);
        \node[font=\bfseries] at (7.55, 4.35) {Hệ thống \texttt{top\_de10\_lite}};
        \node[align=center, font=\scriptsize] at (7.55, 3.55)
            {4.599 / 49.760 LEs (9\%)\\
             2.035 FF (4\%), 0 M9K, 0 DSP};
        \draw[dashed, rounded corners=2pt, fill=gray!8] (5.55, 1.65) rectangle (9.55, 2.65);
        \node[align=center, font=\scriptsize] at (7.55, 2.15)
            {Bộ nhớ Flash nội CFM (Instant-On)\\\texttt{SINGLE IMAGE WITH ERAM}};

        % Khoi ngoai vi ben trai
        \draw[thick, rounded corners=3pt, fill=orange!10] (0.0, 4.7) rectangle (3.6, 5.85);
        \node[align=center] at (1.8, 5.27)
            {\textbf{USB-Blaster II \& 5\,V}\\USB Type-B, JTAG \texttt{.sof}};

        \draw[thick, rounded corners=3pt, fill=green!8] (0.0, 3.3) rectangle (3.6, 4.35);
        \node[align=center] at (1.8, 3.82)
            {\textbf{Dao động 50\,MHz}\\\texttt{PIN\_P11} ($T=20{,}000\text{ ns}$)};

        \draw[thick, rounded corners=3pt, fill=yellow!12] (0.0, 1.8) rectangle (3.6, 2.95);
        \node[align=center] at (1.8, 2.37)
            {\textbf{2 Nút nhấn \texttt{KEY[1:0]}}\\\textit{Schmitt Trigger}, tích cực thấp};

        \draw[thick, rounded corners=3pt, fill=yellow!12] (0.0, 0.25) rectangle (3.6, 1.4);
        \node[align=center] at (1.8, 0.82)
            {\textbf{10 Switch \texttt{SW[9:0]}}\\\texttt{3.3-V LVTTL}, nhập/chế độ};

        % Khoi ngoai vi ben phai
        \draw[thick, rounded corners=3pt, fill=red!8] (11.5, 4.25) rectangle (15.1, 5.55);
        \node[align=center] at (13.3, 4.9)
            {\textbf{3 LED đơn \texttt{LEDR[2:0]}}\\Báo \textit{Busy / Done / Fail}};

        \draw[thick, rounded corners=3pt, fill=red!8] (11.5, 2.25) rectangle (15.1, 3.8);
        \node[align=center] at (13.3, 3.02)
            {\textbf{LED 7 đoạn \texttt{HEX1:HEX0}}\\Anode chung, 8-bit (\texttt{DP=1})};

        \draw[dashed, rounded corners=3pt, fill=gray!8] (11.5, 0.35) rectangle (15.1, 1.7);
        \node[align=center, font=\scriptsize] at (13.3, 1.02)
            {\textbf{Header JP1 (dự phòng)}\\UART tùy chọn \texttt{PIN\_V10, W10}};

        % Ket noi
        \draw[->, thick] (3.6, 5.27) -- node[above, font=\scriptsize]{JTAG} (4.8, 5.27);
        \draw[->, thick] (3.6, 3.82) -- node[above, font=\scriptsize]{Clk} (5.25, 3.82);
        \draw[->, thick] (3.6, 2.37) -- node[above, font=\scriptsize]{2 bit} (5.25, 2.37);
        \draw[->, thick] (3.6, 0.82) -- (4.25, 0.82) |- node[pos=0.72, above, font=\scriptsize]{10 bit} (5.25, 1.9);

        \draw[->, thick] (9.85, 4.3) -- (10.65, 4.3) |- node[pos=0.72, above, font=\scriptsize]{3 bit} (11.5, 4.9);
        \draw[->, thick] (9.85, 3.02) -- node[above, font=\scriptsize]{16 bit} (11.5, 3.02);
        \draw[<->, dashed] (10.3, 1.02) -- (11.5, 1.02);
    \end{tikzpicture}%
    }
    \caption{Sơ đồ khối kết nối ngoại vi trên board Terasic DE10-Lite cho hệ thống RS(15,11)}
    \label{fig:de10-overview}
\end{figure}

\section{Đặc tính kỹ thuật vi mạch FPGA Intel MAX 10 (\texttt{10M50DAF484C7G})}
\label{sec:max10-spec}

Mã linh kiện \texttt{10M50DAF484C7G} quy định cụ thể kiến trúc vi mạch trung tâm trên DE10-Lite: dòng MAX 10 dung lượng 50 nghìn phần tử logic (\texttt{10M50}), biến thể hỗ trợ Flash kép và ADC (\texttt{DA}), vỏ FBGA 484 chân cho tối đa 360 chân I/O (\texttt{F484}), dải nhiệt độ thương mại $0\text{--}85^{\circ}\text{C}$ cấp tốc độ 7 (\texttt{C7}) và chuẩn không chì RoHS6 (\texttt{G}). Theo báo cáo thiết bị của Quartus Prime Fitter (\texttt{ReedSolomon.fit.summary}), vi mạch cung cấp \textbf{49.760 phần tử logic (LEs)} tổ chức trong \textbf{3.110 khối LAB}, \textbf{182 khối RAM M9K} (tổng \textbf{1.677.312 bit}), \textbf{288 bộ nhân $9\times 9$-bit}, \textbf{4 PLL}, \textbf{2 khối ADC} và bộ nhớ Flash cấu hình nội \textbf{CFM/UFM} hỗ trợ khởi động tức thì (\textit{Instant-On} $<10\text{ ms}$) theo chế độ \texttt{SINGLE IMAGE WITH ERAM}.

Trong thiết kế $\mathrm{RS}(15,11)$, hai quyết định ánh xạ phần cứng nổi bật tạo nên đặc trưng tài nguyên riêng của đồ án: (1) phép nhân trên $\mathrm{GF}(256)$ là phép nhân đa thức không mang nhớ modulo $p(u) = u^8 + u^4 + u^3 + u^2 + 1$, do đó không sử dụng các bộ nhân số nguyên $9\times 9$ mà tổng hợp trực tiếp bằng cây cổng XOR trên LE; (2) các bộ đệm dữ liệu 11 byte và 15 byte có dung lượng nhỏ nhưng đòi hỏi truy cập đọc tổ hợp đồng thời và xóa trong 1 chu kỳ nên được ánh xạ hoàn toàn lên flip-flop thay vì M9K. Ngoài ra, tín hiệu \texttt{MAX10\_CLK1\_50} tại \texttt{PIN\_P11} (Bank 3) đi thẳng vào mạng phân phối xung nhịp toàn cục (\textit{Global Clock Network}) để cấp clock đồng bộ cho toàn bộ 2.035 thanh ghi. Bảng~\ref{tab:max10-resources} tổng hợp các thông số tài nguyên danh định của chip.

\begin{table}[!htb]
    \centering
    \caption{Thông số tài nguyên phần cứng của vi mạch Intel MAX 10 (\texttt{10M50DAF484C7G})}
    \label{tab:max10-resources}
    \footnotesize
    \renewcommand{\arraystretch}{1.1}
    \begin{tabular}{p{5.2cm} c p{5.5cm}}
        \toprule
        \textbf{Loại tài nguyên phần cứng} & \textbf{Dung lượng tối đa} & \textbf{Cách khai thác trong đồ án} \\
        \midrule
        Phần tử logic (LEs) / Khối LAB & 49.760 LEs / 3.110 LABs & Thực hiện FSM, mạch $\mathrm{GF}(256)$ và đệm FF (dùng 9\%) \\
        Bộ nhớ khối nhúng (M9K) & 182 khối (1.677.312 bit) & 0\% do đệm 11/15 byte dùng thanh ghi đọc 0 trễ \\
        Bộ nhân nhúng $9\times 9$-bit & 288 phần tử & 0\% do nhân $\mathrm{GF}(256)$ dùng logic XOR tổ hợp \\
        Chân I/O người dùng (FBGA-484) & 360 chân & Dùng 32 chân (13 vào, 19 ra -- chiếm 9\%) \\
        Vòng khóa pha (PLLs) / Flash nội & 4 PLLs / CFM + 1 UFM & Chạy trực tiếp 50\,MHz; chế độ \texttt{SINGLE IMAGE WITH ERAM} \\
        \bottomrule
    \end{tabular}
\end{table}

\section{Hệ thống ngoại vi giao tiếp trên DE10-Lite}
\label{sec:peripherals}

Hệ thống sử dụng 4 nhóm ngoại vi giao tiếp vật lý trên DE10-Lite, được chuẩn hóa cả về cực tính điện áp lẫn lớp đồng bộ chống nhiễu trong RTL:
\begin{itemize}
    \item \textbf{10 công tắc trượt \texttt{SW[9:0]} (tích cực cao, \texttt{3.3-V LVTTL}):} Gạt lên cho mức \texttt{1}, gạt xuống cho mức \texttt{0}, đi qua bộ đồng bộ 2 tầng FF trước khi giải mã theo chế độ. Bit \texttt{SW[8]} chọn pha hoạt động (\texttt{0}: pha nhập liệu; \texttt{1}: pha chạy và xem kết quả). Khi \texttt{SW[8]=0}, \texttt{SW[7:0]} biểu diễn 1 byte dữ liệu đầu vào ($0\text{x00}\dots 0\text{xFF}$). Khi \texttt{SW[8]=1}, \texttt{SW[9]} bật/tắt bộ tiêm lỗi (\texttt{1} = bật), \texttt{SW[7]} chọn số lỗi tiêm vào (\texttt{0} = 2 lỗi trong ngưỡng; \texttt{1} = 3 lỗi vượt ngưỡng), còn \texttt{SW[3:0]} đóng vai trò địa chỉ $0\dots 14$ để đọc từng byte kết quả.
    \item \textbf{2 nút nhấn \texttt{KEY[1:0]} (tích cực thấp, \texttt{3.3 V SCHMITT TRIGGER}):} Bình thường ở mức \texttt{1}, khi nhấn xuống mức \texttt{0}. Mạch \textit{Schmitt Trigger} giúp làm gọn sườn chuyển trạng thái điện áp nhưng vẫn còn nảy phím cơ khí, nên \texttt{KEY[1]} (Xác nhận/Bắt đầu) qua khối lọc dội 10\,ms ($500.000$ chu kỳ ở 50\,MHz) trong \texttt{input\_control}, còn \texttt{KEY[0]} (Reset toàn hệ thống) nối qua \texttt{reset\_sync} để kích hoạt reset bất đồng bộ và nhả reset đồng bộ 2 chu kỳ.
    \item \textbf{3 LED đơn \texttt{LEDR[2:0]} (tích cực cao, \texttt{3.3-V LVTTL}):} \texttt{LEDR[0]} báo \textit{Busy} (sáng khi đang xử lý), \texttt{LEDR[1]} chốt báo \textit{Done} (hoàn tất) và \texttt{LEDR[2]} chốt báo \textit{Fail} (lỗi vượt ngưỡng $T=2$). Tổ hợp \texttt{3'b010} chỉ thị giải mã thành công, \texttt{3'b110} chỉ thị thất bại.
    \item \textbf{2 LED 7 đoạn \texttt{HEX1:HEX0} (Anode chung, tích cực thấp):} Mức \texttt{0} làm thanh LED sáng. Trong mỗi bus 8 bit, \texttt{HEXx[6:0]} nối tới các thanh \texttt{g, f, e, d, c, b, a}, còn \texttt{HEXx[7]} là dấu chấm thập phân \texttt{DP} được module \texttt{hex\_decoder} luôn kéo lên \texttt{1'b1} (tắt hoàn toàn). Ví dụ, giá trị byte \texttt{8'hA7} xuất mã \texttt{HEX1 = 8'h88} (chữ \texttt{A}) và \texttt{HEX0 = 8'hF8} (số \texttt{7}).
\end{itemize}

\section{Quy hoạch và gán chân tín hiệu (Pin Assignment) cho hệ thống}
\label{sec:pin-assignment}

Tệp ràng buộc dự án \texttt{quartus/ReedSolomon.qsf} khai báo thực thể đỉnh \texttt{top\_de10\_lite}, đặt các chân chưa dùng ở chế độ \texttt{AS INPUT TRI-STATED} có điện trở kéo lên yếu và gán toàn bộ \textbf{32 chân tín hiệu} (13 đầu vào, 19 đầu ra) nằm trên các Bank 3, Bank 6 và Bank 7 (đều cấp nguồn I/O $3{,}3\text{ V}$, dòng lái đầu ra 8\,mA). Bảng~\ref{tab:pin-assign} liệt kê đầy đủ toàn bộ 32 chân vật lý trích xuất chính xác từ \texttt{ReedSolomon.qsf} và \texttt{ReedSolomon.fit.rpt} (hai chân \texttt{PIN\_V10}, \texttt{PIN\_W10} trên header \texttt{JP1} được ghi nhận dự phòng cho UART nhưng chưa gán).

\begin{table}[!htb]
    \centering
    \caption{Bảng quy hoạch toàn bộ 32 chân tín hiệu vật lý trên Kit DE10-Lite}
    \label{tab:pin-assign}
    \scriptsize
    \setlength{\tabcolsep}{4pt}
    \renewcommand{\arraystretch}{1.01}
    \begin{tabular}{l c c c l p{4.6cm}}
        \toprule
        \textbf{Tên cổng RTL} & \textbf{Chân FPGA} & \textbf{Hướng} & \textbf{Bank} & \textbf{Chuẩn I/O} & \textbf{Chức năng tín hiệu} \\
        \midrule
        \texttt{MAX10\_CLK1\_50} & \texttt{PIN\_P11} & Input & 3 & \texttt{3.3-V LVTTL} & Xung nhịp hệ thống 50\,MHz \\
        \texttt{KEY[0]} & \texttt{PIN\_B8} & Input & 7 & \texttt{3.3 V SCHMITT TRIGGER} & Reset hệ thống (tích cực thấp) \\
        \texttt{KEY[1]} & \texttt{PIN\_A7} & Input & 7 & \texttt{3.3 V SCHMITT TRIGGER} & Xác nhận nạp / Bắt đầu chạy \\
        \texttt{SW[0]} & \texttt{PIN\_C10} & Input & 7 & \texttt{3.3-V LVTTL} & Bit dữ liệu 0 / Địa chỉ đọc bit 0 \\
        \texttt{SW[1]} & \texttt{PIN\_C11} & Input & 7 & \texttt{3.3-V LVTTL} & Bit dữ liệu 1 / Địa chỉ đọc bit 1 \\
        \texttt{SW[2]} & \texttt{PIN\_D12} & Input & 7 & \texttt{3.3-V LVTTL} & Bit dữ liệu 2 / Địa chỉ đọc bit 2 \\
        \texttt{SW[3]} & \texttt{PIN\_C12} & Input & 7 & \texttt{3.3-V LVTTL} & Bit dữ liệu 3 / Địa chỉ đọc bit 3 \\
        \texttt{SW[4]} & \texttt{PIN\_A12} & Input & 7 & \texttt{3.3-V LVTTL} & Bit dữ liệu 4 (khi \texttt{SW[8]=0}) \\
        \texttt{SW[5]} & \texttt{PIN\_B12} & Input & 7 & \texttt{3.3-V LVTTL} & Bit dữ liệu 5 (khi \texttt{SW[8]=0}) \\
        \texttt{SW[6]} & \texttt{PIN\_A13} & Input & 7 & \texttt{3.3-V LVTTL} & Bit dữ liệu 6 (khi \texttt{SW[8]=0}) \\
        \texttt{SW[7]} & \texttt{PIN\_A14} & Input & 7 & \texttt{3.3-V LVTTL} & Bit dữ liệu 7 / Chọn 2 (\texttt{0}) hay 3 lỗi (\texttt{1}) \\
        \texttt{SW[8]} & \texttt{PIN\_B14} & Input & 7 & \texttt{3.3-V LVTTL} & Chọn chế độ: \texttt{0} = Nhập, \texttt{1} = Chạy/Xem \\
        \texttt{SW[9]} & \texttt{PIN\_F15} & Input & 7 & \texttt{3.3-V LVTTL} & Bật/tắt tiêm lỗi (khi \texttt{SW[8]=1}) \\
        \midrule
        \texttt{LEDR[0]} & \texttt{PIN\_A8} & Output & 7 & \texttt{3.3-V LVTTL} & Đèn báo \textit{Busy} (đang xử lý) \\
        \texttt{LEDR[1]} & \texttt{PIN\_A9} & Output & 7 & \texttt{3.3-V LVTTL} & Đèn báo \textit{Done} (đã hoàn tất) \\
        \texttt{LEDR[2]} & \texttt{PIN\_A10} & Output & 7 & \texttt{3.3-V LVTTL} & Đèn báo \textit{Fail} (vượt ngưỡng $T=2$) \\
        \texttt{HEX0[0]} & \texttt{PIN\_C14} & Output & 7 & \texttt{3.3-V LVTTL} & Thanh \texttt{a} của \texttt{HEX0} (nibble thấp) \\
        \texttt{HEX0[1]} & \texttt{PIN\_E15} & Output & 7 & \texttt{3.3-V LVTTL} & Thanh \texttt{b} của \texttt{HEX0} \\
        \texttt{HEX0[2]} & \texttt{PIN\_C15} & Output & 7 & \texttt{3.3-V LVTTL} & Thanh \texttt{c} của \texttt{HEX0} \\
        \texttt{HEX0[3]} & \texttt{PIN\_C16} & Output & 7 & \texttt{3.3-V LVTTL} & Thanh \texttt{d} của \texttt{HEX0} \\
        \texttt{HEX0[4]} & \texttt{PIN\_E16} & Output & 7 & \texttt{3.3-V LVTTL} & Thanh \texttt{e} của \texttt{HEX0} \\
        \texttt{HEX0[5]} & \texttt{PIN\_D17} & Output & 7 & \texttt{3.3-V LVTTL} & Thanh \texttt{f} của \texttt{HEX0} \\
        \texttt{HEX0[6]} & \texttt{PIN\_C17} & Output & 7 & \texttt{3.3-V LVTTL} & Thanh \texttt{g} của \texttt{HEX0} \\
        \texttt{HEX0[7]} & \texttt{PIN\_D15} & Output & 7 & \texttt{3.3-V LVTTL} & Dấu chấm \texttt{DP} của \texttt{HEX0} (luôn tắt) \\
        \texttt{HEX1[0]} & \texttt{PIN\_C18} & Output & 7 & \texttt{3.3-V LVTTL} & Thanh \texttt{a} của \texttt{HEX1} (nibble cao) \\
        \texttt{HEX1[1]} & \texttt{PIN\_D18} & Output & 6 & \texttt{3.3-V LVTTL} & Thanh \texttt{b} của \texttt{HEX1} \\
        \texttt{HEX1[2]} & \texttt{PIN\_E18} & Output & 6 & \texttt{3.3-V LVTTL} & Thanh \texttt{c} của \texttt{HEX1} \\
        \texttt{HEX1[3]} & \texttt{PIN\_B16} & Output & 7 & \texttt{3.3-V LVTTL} & Thanh \texttt{d} của \texttt{HEX1} \\
        \texttt{HEX1[4]} & \texttt{PIN\_A17} & Output & 7 & \texttt{3.3-V LVTTL} & Thanh \texttt{e} của \texttt{HEX1} \\
        \texttt{HEX1[5]} & \texttt{PIN\_A18} & Output & 7 & \texttt{3.3-V LVTTL} & Thanh \texttt{f} của \texttt{HEX1} \\
        \texttt{HEX1[6]} & \texttt{PIN\_B17} & Output & 7 & \texttt{3.3-V LVTTL} & Thanh \texttt{g} của \texttt{HEX1} \\
        \texttt{HEX1[7]} & \texttt{PIN\_A16} & Output & 7 & \texttt{3.3-V LVTTL} & Dấu chấm \texttt{DP} của \texttt{HEX1} (luôn tắt) \\
        \bottomrule
    \end{tabular}
\end{table}

Để kiểm tra hợp chuẩn định thời trên TimeQuest, tệp \texttt{quartus/timing.sdc} định nghĩa nguồn xung nhịp chính và độ bất định xung nhịp như sau:
\begin{lstlisting}[numbers=none]
create_clock -name sys_clk -period 20.000 [get_ports {MAX10_CLK1_50}]
derive_clock_uncertainty
\end{lstlisting}
Chu kỳ \texttt{20.000 ns} tương ứng tần số mục tiêu 50\,MHz, trong khi \texttt{derive\_clock\_uncertainty} tự động nội suy biên độ jitter/skew giữa các thanh ghi MAX 10. Các chân \texttt{KEY}, \texttt{SW}, \texttt{LEDR} và \texttt{HEX} là giao tiếp người dùng bất đồng bộ tốc độ chậm (được đồng bộ nội bộ bằng chuỗi 2 FF) nên không áp đặt ràng buộc \texttt{set\_input\_delay} hay \texttt{set\_output\_delay} giả định.


% ============================================================================
% CHUONG 3
% ============================================================================
\chapter{Thực hiện Reed-Solomon trên FPGA}
\label{ch:rs-fpga}

Chương này phân tích chi tiết kiến trúc phần cứng RTL của bộ mã hóa -- giải mã $\mathrm{RS}(15,11)$ trên FPGA, bao gồm tổ chức hệ thống 21 tệp Verilog, giao thức bắt tay luồng, thiết kế các khối số học $\mathrm{GF}(256)$, bộ mã hóa LFSR, 6 module thành phần của bộ giải mã, các khối giao tiếp bo mạch cùng kết quả mô phỏng và tổng hợp Quartus Prime thực tế.

\section{Kiến trúc tổng thể hệ thống trên FPGA}
\label{sec:top-arch}

Hệ thống phần cứng được xây dựng theo cấu trúc phân cấp gồm tệp tham số \texttt{rtl/include/rs\_params.vh} và \textbf{21 tệp Verilog RTL} chia thành 7 phân hệ: \texttt{gf\_arithmetic}, \texttt{rs\_encoder}, \texttt{rs\_decoder}, \texttt{memory}, \texttt{controller}, \texttt{demo} và \texttt{board}. Các hằng số trong \texttt{rs\_params.vh} ấn định: $m = 8$, \texttt{GF\_PRIM\_POLY = 9'h11D} ($u^8+u^4+u^3+u^2+1$), \texttt{GF\_ALPHA = 8'h02}, $N = 15$, $K = 11$, $R = 4$, $T = 2$ (rút gọn $240$ symbol đầu từ mã mẹ $\mathrm{RS}(255,251)$), nghiệm liên tiếp bắt đầu tại $F = 0$ với bước $1$. Hình~\ref{fig:top-arch} biểu diễn sơ đồ khối kết nối toàn bộ các thực thể bên trong \texttt{top\_de10\_lite}.

\begin{figure}[!htb]
    \centering
    \resizebox{\textwidth}{!}{%
    \begin{tikzpicture}[>=stealth, font=\footnotesize]
        % Tang giao tiep ngo vao
        \draw[thick, rounded corners=2pt, fill=blue!6] (0.0, 6.2) rectangle (3.7, 7.25);
        \node[align=center] at (1.85, 6.72) {\texttt{reset\_sync} (\texttt{u\_reset\_sync})\\Async assert, 2-FF sync};

        \draw[thick, rounded corners=2pt, fill=blue!6] (4.7, 6.2) rectangle (10.7, 7.25);
        \node[align=center] at (7.7, 6.72) {\texttt{input\_control} (\texttt{u\_input\_ctrl})\\2-FF sync \texttt{KEY[1], SW[9:0]}; Debounce 10\,ms};

        \node[font=\scriptsize\bfseries] at (1.85, 7.75) {\texttt{MAX10\_CLK1\_50}, \texttt{KEY[0]}};
        \draw[->, thick] (1.85, 7.55) -- (1.85, 7.25);
        \node[font=\scriptsize\bfseries] at (7.7, 7.75) {\texttt{KEY[1]}, \texttt{SW[9:0]}};
        \draw[->, thick] (7.7, 7.55) -- (7.7, 7.25);

        % Duong du lieu streaming giua
        \draw[thick, rounded corners=2pt, fill=green!7] (0.0, 3.8) rectangle (2.6, 5.1);
        \node[align=center] at (1.3, 4.45) {\texttt{data\_buffer}\\\texttt{u\_buf\_input}\\($11\times 8$-bit)};

        \draw[thick, rounded corners=2pt, fill=orange!9] (3.3, 3.8) rectangle (5.9, 5.1);
        \node[align=center] at (4.6, 4.45) {\texttt{rs\_encoder}\\\texttt{u\_encoder}\\LFSR $R=4$};

        \draw[thick, rounded corners=2pt, fill=red!8] (6.6, 3.8) rectangle (9.2, 5.1);
        \node[align=center] at (7.9, 4.45) {\texttt{error\_injector}\\\texttt{u\_injector}\\0 / 2 / 3 lỗi};

        \draw[thick, rounded corners=2pt, fill=orange!11] (9.9, 3.8) rectangle (12.6, 5.1);
        \node[align=center] at (11.25, 4.45) {\texttt{rs\_decoder}\\\texttt{u\_decoder}\\6 khối + Re-verify};

        \draw[thick, rounded corners=2pt, fill=green!7] (13.3, 3.8) rectangle (15.9, 5.1);
        \node[align=center] at (14.6, 4.45) {\texttt{data\_buffer}\\\texttt{u\_buf\_output}\\($15\times 8$-bit)};

        \draw[->, thick] (2.6, 4.45) -- node[above, font=\scriptsize]{\texttt{v/r}} (3.3, 4.45);
        \draw[->, thick] (5.9, 4.45) -- node[above, font=\scriptsize]{\texttt{v/r/l}} (6.6, 4.45);
        \draw[->, thick] (9.2, 4.45) -- node[above, font=\scriptsize]{\texttt{v/r/l}} (9.9, 4.45);
        \draw[->, thick] (12.6, 4.45) -- node[above, font=\scriptsize]{\texttt{v/r/l}} (13.3, 4.45);

        % Hang duoi
        \draw[thick, rounded corners=2pt, fill=gray!10] (0.0, 0.7) rectangle (3.3, 2.15);
        \node[align=center] at (1.65, 1.42) {\texttt{result\_checker}\\\texttt{u\_checker}\\So sánh 11 byte};

        \draw[thick, rounded corners=2pt, fill=purple!8] (4.3, 0.7) rectangle (9.4, 2.15);
        \node[align=center] at (6.85, 1.42) {\texttt{rs\_controller} (\texttt{u\_controller})\\FSM 4 trạng thái (\texttt{S\_IDLE..S\_DONE})};

        \draw[thick, rounded corners=2pt, fill=gray!10] (10.1, 0.7) rectangle (12.4, 2.15);
        \node[align=center] at (11.25, 1.42) {\texttt{cycle\_counter}\\\texttt{u\_counter}\\Đếm 32-bit};

        \draw[thick, rounded corners=2pt, fill=blue!7] (13.1, 0.7) rectangle (15.9, 2.15);
        \node[align=center] at (14.5, 1.42) {\texttt{display\_control}\\\& \texttt{hex\_decoder}};

        % Ket noi duong dieu khien (khong chong cheo nhan)
        \draw[->, thick] (6.1, 6.2) -- (6.1, 2.15);
        \node[font=\scriptsize, fill=white, inner sep=1pt, anchor=east] at (5.95, 5.6) {\texttt{confirm\_pulse, sw\_sync, sw\_latched}};

        \draw[->, thick] (1.1, 3.8) -- (1.1, 2.15);
        \draw[->, thick] (14.6, 3.8) -- (14.6, 2.75) -| (2.5, 2.15);
        \draw[->, thick] (14.0, 2.75) -| (8.2, 2.15);
        \node[font=\scriptsize, fill=white, inner sep=1pt] at (11.4, 2.75) {\texttt{rd\_data / chk\_data}};

        \draw[->, thick] (6.85, 2.15) -- (6.85, 3.25) -| (4.6, 3.8);
        \draw[->, thick] (6.85, 3.25) -| (7.9, 3.8);
        \draw[->, thick] (6.85, 3.25) -| (11.25, 3.8);
        \draw[->, thick] (6.85, 3.25) -| (1.9, 3.8);

        \draw[->, thick] (9.4, 1.7) -- (10.1, 1.7);
        \draw[->, thick] (4.3, 1.42) -- (3.3, 1.42);
        \draw[->, thick] (9.4, 1.1) -- node[below, font=\scriptsize]{\texttt{byte, status}} (13.1, 1.1);

        \node[font=\scriptsize\bfseries] at (14.5, 0.05) {\texttt{LEDR[2:0]}, \texttt{HEX1[7:0]}, \texttt{HEX0[7:0]}};
        \draw[->, thick] (14.5, 0.7) -- (14.5, 0.25);
    \end{tikzpicture}%
    }
    \caption{Sơ đồ khối kiến trúc tổng thể hệ thống \texttt{top\_de10\_lite} trên FPGA}
    \label{fig:top-arch}
\end{figure}

\subsection*{Giao thức luồng (\textit{Valid/Ready}), bộ đệm \texttt{data\_buffer} và điều phối \texttt{rs\_controller}}
Chuỗi truyền dữ liệu nối tiếp giữa \texttt{u\_buf\_input}, \texttt{u\_encoder}, \texttt{u\_injector}, \texttt{u\_decoder} và \texttt{u\_buf\_output} tuân thủ giao thức bắt tay \texttt{valid}/\texttt{ready} có \textit{backpressure}. Symbol 8 bit chỉ dịch chuyển khi thỏa mãn $\text{fire} = \texttt{valid} \wedge \texttt{ready}$. Thanh ghi đầu ra chỉ nhận byte mới khi $\texttt{out\_free} = \texttt{!out\_valid} \vee \texttt{out\_ready}$, ngăn chặn ghi đè khi tầng sau hạ \texttt{ready = 0}. Tín hiệu \texttt{out\_last} lên cao tại byte thứ 15 ($N-1 = 14$), \texttt{busy} giữ mức cao suốt phiên chạy và \texttt{done} phát xung 1 chu kỳ khi hoàn tất.

Module \texttt{data\_buffer} triển khai mảng thanh ghi \texttt{reg [7:0] mem [0:DEPTH-1]} hỗ trợ đồng thời ghi ngẫu nhiên (\texttt{wr\_en/wr\_addr/wr\_data}), hai cổng đọc tổ hợp 0 chu kỳ trễ (\texttt{rd\_addr} cho hiển thị và \texttt{chk\_addr} cho kiểm tra) cùng giao diện truyền/nhận luồng (\texttt{stream\_start}). Thiết kế khởi tạo hai bản sao: \texttt{u\_buf\_input} (\texttt{DEPTH=11}) và \texttt{u\_buf\_output} (\texttt{DEPTH=15}).

Bộ điều phối \texttt{rs\_controller} sử dụng máy trạng thái 4 trạng thái (\texttt{S\_IDLE}, \texttt{S\_START}, \texttt{S\_RUN}, \texttt{S\_DONE}):
\begin{itemize}
    \item \textbf{Pha nhập liệu (\texttt{sw\_sync[8] == 0}):} Mỗi xung \texttt{confirm\_pulse} ghi \texttt{sw\_latched[7:0]} vào \texttt{u\_buf\_input} tại con trỏ \texttt{in\_cnt} ($0\dots 10$ quay vòng). Đầu ra \texttt{byte\_to\_display} nối trực tiếp với \texttt{sw\_sync[7:0]} và bật \texttt{clear\_status = 1}. Sườn xuống của \texttt{SW[8]} (\texttt{1->0}) tự động đặt lại \texttt{in\_cnt <= 0} và \texttt{has\_run <= 0}.
    \item \textbf{Pha chạy và xem kết quả (\texttt{sw\_sync[8] == 1}):} Tại \texttt{S\_IDLE}, xung \texttt{confirm\_pulse} bật \texttt{busy <= 1}, phát \texttt{clear\_status\_pulse} và chuyển sang \texttt{S\_START}. Ở \texttt{S\_START}, hệ thống phát đồng thời xung 1 chu kỳ tới \texttt{buf\_out\_clear}, \texttt{counter\_clear}, \texttt{counter\_start}, \texttt{encoder\_start}, \texttt{injector\_start}, \texttt{decoder\_start} và \texttt{buf\_in\_stream\_start}, rồi sang \texttt{S\_RUN} chờ \texttt{decoder\_done}. Khi giải mã xong, \texttt{rs\_controller} chốt \texttt{fail <= decoder\_fail}, phát \texttt{done}, \texttt{counter\_stop}, \texttt{checker\_start}, bật \texttt{has\_run <= 1} và qua \texttt{S\_DONE} trở về \texttt{S\_IDLE}. Khi xem lại kết quả (\texttt{has\_run = 1}), \texttt{byte\_to\_display} đọc \texttt{u\_buf\_output} tại \texttt{sw\_sync[3:0]} (giữ về 0 nếu giá trị $\ge 15$).
\end{itemize}

\section{Thiết kế RTL khối Số học $\mathrm{GF}(256)$ và Bộ mã hóa (\texttt{rs\_encoder})}
\label{sec:gf-encoder}

Ba khối trong \texttt{rtl/gf\_arithmetic/} thực hiện toàn bộ phép toán trường $\mathrm{GF}(256)$ theo dạng mạch tổ hợp thuần túy (0 chu kỳ trễ):
\begin{itemize}
    \item \textbf{\texttt{gf256\_add}:} Thực hiện phép cộng không mang nhớ qua 8 cổng XOR song song (\texttt{assign y = a \textasciicircum{} b;}).
    \item \textbf{\texttt{gf256\_mul}:} Tính tích $(a \cdot b) \bmod p(u)$ qua vòng lặp mở 8 tầng dịch trái và rút gọn theo \texttt{POLY\_LOW = 8'h1D} (\textit{Russian peasant shift-and-reduce}). So với bảng tra log/antilog kép trên RAM, cấu trúc tổ hợp này không tốn bộ nhớ M9K, không có độ trễ đọc và cho phép Quartus rút gọn tối đa nhánh hằng khi nhân với các hệ số cố định trong bộ mã hóa hay khối syndrome.
    \item \textbf{\texttt{gf256\_inv}:} Tính $a^{-1} = a^{254}$ qua bảng tra 256 mục \texttt{INV\_LUT} sinh tự động lúc \textit{elaboration} bởi hàm \texttt{build\_inv\_lut}. Khi $a = \text{0x00}$, khối xuất \texttt{y = 8'h00} và bật \texttt{zero\_err = 1} (\texttt{valid = 0}) để phục vụ bắt lỗi chia cho 0 trong Forney.
\end{itemize}

Bộ mã hóa \texttt{rs\_encoder} tạo từ mã hệ thống $c(x) = x^R m(x) + \big(x^R m(x) \bmod g(x)\big)$ dựa trên đa thức sinh $g(x) = x^4 + \text{0x0F}x^3 + \text{0x36}x^2 + \text{0x78}x + \text{0x40}$ (tính tự động trong tham số 40-bit \texttt{GEN}). Kiến trúc sử dụng LFSR 4 thanh ghi \texttt{parity[0:3]} (\texttt{parity[0]} chứa hệ số bậc $x^3$) và FSM 4 trạng thái (\texttt{S\_IDLE}, \texttt{S\_DATA}, \texttt{S\_PARITY}, \texttt{S\_DRAIN}): trong $K=11$ chu kỳ ở \texttt{S\_DATA}, từng byte dữ liệu vừa đi thẳng ra \texttt{data\_out}, vừa đi vào mạch phản hồi \texttt{feedback = data\_in \textasciicircum{} parity[0]} nhân với các hệ số \texttt{GEN[8*(R-1-g) +: 8]} để cập nhật \texttt{parity[0..3]}; ở \texttt{S\_PARITY}, 4 byte kiểm tra được đẩy tuần tự ra ngoài với \texttt{out\_last = 1} tại byte thứ 15, rồi chuyển sang \texttt{S\_DRAIN} đợi xác nhận \texttt{out\_fire} trước khi phát xung \texttt{done}. Đoạn mã Verilog sau trích dẫn khối tạo nhánh phản hồi LFSR của \texttt{rs\_encoder}:

\begin{lstlisting}[language=Verilog]
gf256_add u_fb_add ( .a(data_in), .b(parity[0]), .y(feedback) );

genvar g;
generate
    for (g = 0; g < R; g = g + 1) begin : gen_lfsr_taps
        gf256_mul u_tap_mul (
            .a (feedback),
            .b (GEN[8*(R-1-g) +: 8]),
            .y (fb_mul[g])
        );
        if (g < R - 1) begin : gen_tap_add
            gf256_add u_tap_add (
                .a (parity[g+1]), .b (fb_mul[g]), .y (par_next[g])
            );
        end else begin : gen_tap_last
            assign par_next[g] = fb_mul[g];
        end
    end
endgenerate
\end{lstlisting}

\section{Thiết kế RTL Bộ giải mã (\texttt{rs\_decoder}) và 6 Module thành phần}
\label{sec:rs-decoder}

Bộ giải mã \texttt{rs\_decoder} sử dụng giao diện luồng đồng nhất với bộ mã hóa (bổ sung cờ \texttt{fail}) và điều phối 6 module toán học thông qua máy trạng thái 12 trạng thái \texttt{state[3:0]} (Hình~\ref{fig:decoder-fsm}).

\begin{figure}[!htb]
    \centering
    \resizebox{0.98\textwidth}{!}{%
    \begin{tikzpicture}[>=stealth, font=\footnotesize]
        % Cot trai
        \node[draw, thick, rounded corners=3pt, fill=blue!10, minimum width=2.4cm, minimum height=0.65cm] (Sidle) at (2.0, 7.6) {\texttt{S\_IDLE}};
        \node[draw, thick, rounded corners=3pt, fill=green!10, minimum width=2.4cm, minimum height=0.65cm] (Srecv) at (2.0, 6.2) {\texttt{S\_RECV} (16)};
        \node[draw, thick, rounded corners=3pt, fill=yellow!15, minimum width=2.7cm, minimum height=0.65cm] (Swait1) at (2.0, 4.8) {\texttt{S\_WAIT\_SYN1} (1)};
        \node[draw, thick, rounded corners=3pt, fill=orange!12, minimum width=2.7cm, minimum height=0.65cm] (Sbm) at (2.0, 3.3) {\texttt{S\_BM} (11)};
        \node[draw, thick, rounded corners=3pt, fill=orange!12, minimum width=2.7cm, minimum height=0.65cm] (Somega) at (2.0, 1.9) {\texttt{S\_OMEGA} (6)};
        \node[draw, thick, rounded corners=3pt, fill=orange!12, minimum width=2.7cm, minimum height=0.65cm] (Schien) at (2.0, 0.5) {\texttt{S\_CHIEN} (17)};

        % Cot giua
        \node[draw, thick, rounded corners=3pt, fill=orange!12, minimum width=2.8cm, minimum height=0.65cm] (Sforney) at (7.0, 0.5) {\texttt{S\_FORNEY} (17)};
        \node[draw, thick, rounded corners=3pt, fill=orange!12, minimum width=2.8cm, minimum height=0.65cm] (Scorr) at (7.0, 1.9) {\texttt{S\_CORRECT} (17)};
        \node[draw, thick, rounded corners=3pt, fill=purple!10, minimum width=3.1cm, minimum height=0.65cm] (Svfeed) at (7.0, 3.3) {\texttt{S\_VERIFY\_FEED} (16)};
        \node[draw, thick, rounded corners=3pt, fill=purple!10, minimum width=3.1cm, minimum height=0.65cm] (Svwait) at (7.0, 4.8) {\texttt{S\_VERIFY\_WAIT} (1)};

        % Cot phai
        \node[draw, thick, rounded corners=3pt, fill=green!12, minimum width=2.8cm, minimum height=0.65cm] (Soprep) at (12.8, 4.8) {\texttt{S\_OUT\_PREP} (1)};
        \node[draw, thick, rounded corners=3pt, fill=green!12, minimum width=2.8cm, minimum height=0.65cm] (Sout) at (12.8, 6.8) {\texttt{S\_OUT} (15)};

        % Chuyen trang thai
        \draw[->, thick] (Sidle) -- node[right, font=\scriptsize]{\texttt{start=1}} (Srecv);
        \draw[->, thick] (Srecv) -- node[right, font=\scriptsize]{15 byte} (Swait1);
        \draw[->, thick] (Swait1) -- node[right, font=\scriptsize]{\texttt{syn\_zero=0}} (Sbm);
        \draw[->, thick] (Sbm) -- node[right, font=\scriptsize]{hợp lệ} (Somega);
        \draw[->, thick] (Somega) -- node[right, font=\scriptsize]{\texttt{ee\_done}} (Schien);
        \draw[->, thick] (Schien) -- node[above, font=\scriptsize]{đủ nghiệm} (Sforney);
        \draw[->, thick] (Sforney) -- node[right, font=\scriptsize]{hợp lệ} (Scorr);
        \draw[->, thick] (Scorr) -- node[right, font=\scriptsize]{\texttt{ec\_done}} (Svfeed);
        \draw[->, thick] (Svfeed) -- node[right, font=\scriptsize]{15 byte} (Svwait);
        \draw[->, thick] (Svwait) -- node[above, font=\scriptsize, align=center]{\texttt{syn\_zero=1/0}\\\texttt{fail=0/1}} (Soprep);
        \draw[->, thick] (Soprep) -- (Sout);
        \draw[->, thick] (Sout) |- node[pos=0.75, above, font=\scriptsize]{\texttt{out\_fire \& out\_cnt==14 (done=1)}} (Sidle);

        % Fast-path
        \draw[->, thick, blue!75!black] (3.35, 4.95) -- (4.35, 4.95) -- (4.35, 5.85) -| (12.8, 5.13);
        \node[font=\scriptsize\bfseries, blue!75!black, fill=white, inner sep=1pt] at (8.4, 6.1)
            {Fast-path (\texttt{syn\_zero=1}): \texttt{corr\_buf <= rx\_buf, fail <= 0}};

        % Early Abort
        \draw[->, dashed, red!75!black] (3.35, 3.3) -- (4.35, 3.3) |- (4.35, -0.3) -| (12.8, 4.47);
        \draw[->, dashed, red!75!black] (Schien.south) |- (4.35, -0.3);
        \draw[->, dashed, red!75!black] (Sforney.south) -- (7.0, -0.3);
        \node[font=\scriptsize, red!75!black, fill=white, inner sep=1pt] at (7.5, -0.65)
            {Early Abort (\texttt{bm\_fail} / \texttt{chien\_fail} / \texttt{forney\_fail}): \texttt{fail <= 1, corr\_buf <= rx\_buf}};
    \end{tikzpicture}%
    }
    \caption{Lưu đồ máy trạng thái 12 trạng thái và hai luồng xử lý của \texttt{rs\_decoder}}
    \label{fig:decoder-fsm}
\end{figure}

Kiến trúc \texttt{rs\_decoder} nổi bật với ba cơ chế tối ưu và bảo vệ an toàn giải mã:
\begin{enumerate}
    \item \textbf{Nhánh tắt không lỗi (\textit{Fast-path}):} Trong khi nhận 15 byte ở \texttt{S\_RECV}, dữ liệu đồng thời nạp vào \texttt{rx\_buf} và \texttt{u\_syndrome}. Nếu ở \texttt{S\_WAIT\_SYN1} có \texttt{syn\_zero\_out = 1}, mạch bỏ qua 5 khối giải mã phía sau và nhảy thẳng tới \texttt{S\_OUT\_PREP}, giảm trễ tính toán giữa byte vào cuối và \texttt{out\_valid} xuống còn đúng \textbf{2 chu kỳ} (tổng thời gian 1 khối sạch là \textbf{34 chu kỳ}).
    \item \textbf{Tái sử dụng \texttt{u\_syndrome} cho Re-verify hai lớp:} Thông qua bộ dồn kênh theo tín hiệu \texttt{state == S\_VERIFY\_FEED}, cùng một thực thể \texttt{u\_syndrome} được dùng lại để tính lại hội chứng cho từ mã ứng viên \texttt{candidate\_buf} sau bước \texttt{error\_corrector}. Chỉ khi syndrome lần 2 bằng 0 thì \texttt{fail} mới bằng 0 và xuất \texttt{candidate\_buf}; nếu khác 0 (trường hợp giải mã sai do $>T$ lỗi), mạch bật \texttt{fail <= 1} và khôi phục bản nhận gốc \texttt{rx\_buf}.
    \item \textbf{Thoát sớm (\textit{Early Abort}):} Mọi cờ \texttt{fail} phát sinh từ BM, Chien Search hoặc Forney lập tức đưa FSM sang \texttt{S\_OUT\_PREP} với \texttt{fail <= 1} để tiết kiệm chu kỳ.
\end{enumerate}

Hoạt động chi tiết của 6 khối thành phần bên trong \texttt{rs\_decoder} (tổng hợp chu kỳ xử lý trong Bảng~\ref{tab:decoder-submodules}) như sau:
\begin{itemize}
    \item \textbf{\texttt{syndrome.v}:} FSM 2 trạng thái (\texttt{S\_IDLE}, \texttt{S\_RECV}) gồm 4 nhánh Horner song song ứng với $\beta_j = \alpha^j \in \{0\text{x01}, 0\text{x02}, 0\text{x04}, 0\text{x08}\}$, cập nhật $\text{acc}[j] \leftarrow (\text{acc}[j]\cdot \beta_j) \oplus \texttt{data\_in}$ trong 15 chu kỳ nhận và xuất 32-bit \texttt{syn} kèm cờ \texttt{syn\_zero}.
    \item \textbf{\texttt{berlekamp\_massey.v}:} FSM 4 trạng thái (\texttt{S\_IDLE}, \texttt{S\_DISC}, \texttt{S\_UPDATE}, \texttt{S\_FINISH}) lặp qua $q=0\dots 3$ (9 chu kỳ lõi, 11 chu kỳ wrapper). Trong \texttt{S\_DISC}, mạch tính độ lệch $d = S_q \oplus \bigoplus_{j=1}^L C_j S_{q-j}$; trong \texttt{S\_UPDATE}, dùng \texttt{gf256\_inv(b)} tính $\text{scale} = d \cdot b^{-1}$ để cập nhật $C(z)$, $B(z)$ và bậc $L$. Tại \texttt{S\_FINISH}, điều kiện \texttt{fail <= (!is\_syn\_all\_zero) \&\& (L == 0 || L > T)} chặn các đa thức vượt bậc $T=2$.
    \item \textbf{\texttt{error\_evaluator.v}:} Tính tuần tự 4 hệ số $\Omega_q = \bigoplus_{j=0}^{\min(q,L)} \lambda_j S_{q-j}$ ($q=0\dots 3$) trong 4 chu kỳ lõi (6 chu kỳ wrapper) bằng mảng 4 bộ \texttt{gf256\_mul}.
    \item \textbf{\texttt{chien\_search.v}:} FSM 2 trạng thái (\texttt{S\_IDLE}, \texttt{S\_RUN}) khởi tạo $z_0 = \alpha^{-14} = \text{0x58}$ và nhân bước $\alpha = \text{0x02}$ để kiểm tra $\Lambda(z_i)=0$ tại đúng 15 vị trí thực tế $i=0\dots 14$ trong 15 chu kỳ (17 chu kỳ wrapper). Do mã rút gọn 240 vị trí đầu, nếu $\Lambda(z)$ có nghiệm rơi vào vùng rút gọn hoặc nghiệm bội thì số nghiệm đếm được không khớp bậc $L$ và điều kiện \texttt{fail <= (found\_next != deg)} sẽ báo lỗi.
    \item \textbf{\texttt{forney.v}:} Chạy song song $X_i$ (từ $\alpha^{14} = \text{0x13}$, bước $\alpha^{-1}=\text{0x8E}$) và $z_i$ (từ $\text{0x58}$, bước $\alpha$) trong 15 chu kỳ (17 chu kỳ wrapper) để tính $E_i = X_i \cdot \Omega(z_i)/\Lambda'(z_i)$ với đạo hàm bậc lẻ $\Lambda'(z_i) = \lambda_1 \oplus \lambda_3 z_i^2$. Nếu gặp $\Lambda'(z_i)=0$ tại vị trí có lỗi (\texttt{err\_mask[i]=1}), cờ \texttt{fail} lập tức kích hoạt.
    \item \textbf{\texttt{error\_corrector.v}:} Thực hiện phép XOR tuần tự $\texttt{corrected[i]} \leftarrow \texttt{rx[i]} \oplus \texttt{errors[i]}$ cho 15 byte (120 bit) trong 15 chu kỳ (17 chu kỳ wrapper).
\end{itemize}

\begin{table}[!htb]
    \centering
    \caption{Tổng hợp kiến trúc và số chu kỳ xử lý của 6 module giải mã trong \texttt{rs\_decoder}}
    \label{tab:decoder-submodules}
    \footnotesize
    \renewcommand{\arraystretch}{1.1}
    \begin{tabular}{l c c p{6.2cm}}
        \toprule
        \textbf{Module thành phần} & \textbf{Chu kỳ lõi} & \textbf{Chu kỳ FSM} & \textbf{Đặc điểm phần cứng \& cơ chế an toàn} \\
        \midrule
        \texttt{syndrome} & 15 & 16 / 16 & 4 nhánh Horner ($\beta_j=\alpha^j$); dùng chung pha nhận \& Re-verify \\
        \texttt{berlekamp\_massey} & 9 & 11 & Vòng lặp $q=0..3$; \texttt{fail} khi $L=0$ hoặc $L>2$ \\
        \texttt{error\_evaluator} & 4 & 6 & Chập vòng $q=0..3$ tính $\Omega(z)$ \\
        \texttt{chien\_search} & 15 & 17 & Quét $z=\alpha^{-14}\dots 1$; \texttt{fail} khi số nghiệm $\neq L$ \\
        \texttt{forney} & 15 & 17 & Tính $X_i \Omega(z_i)/\Lambda'(z_i)$; \texttt{fail} khi $\Lambda'(z_i)=0$ \\
        \texttt{error\_corrector} & 15 & 17 & XOR sửa tuần tự 15 byte (120 bit) \\
        Đợi, kiểm và xuất luồng & -- & 19 (18 FP) & Gồm \texttt{S\_IDLE}(1), \texttt{S\_WAIT\_SYN1}(1), \texttt{S\_VERIFY\_WAIT}(1), \texttt{S\_OUT\_PREP}(1), \texttt{S\_OUT}(15) \\
        \midrule
        \textbf{Tổng (Fast-path / Đủ chuỗi)} & -- & \textbf{34 / 119} & Tính từ lúc \texttt{start=1} ở \texttt{S\_IDLE} tới xung \texttt{done=1} \\
        \bottomrule
    \end{tabular}
\end{table}

\section{Thiết kế các khối phụ trợ Demo và Giao tiếp ngoại vi}
\label{sec:demo-board}

Các khối trong \texttt{rtl/demo/} và \texttt{rtl/board/} hoàn thiện vòng khép kín trên DE10-Lite:
\begin{itemize}
    \item \textbf{\texttt{error\_injector.v}:} FSM 3 trạng thái (\texttt{S\_IDLE}, \texttt{S\_PASS}, \texttt{S\_DRAIN}) đặt giữa \texttt{u\_encoder} và \texttt{u\_decoder}. Khi bắt đầu, khối chốt \texttt{enable = SW[9]}, \texttt{two\_errors = 1'b1} và \texttt{three\_errors = SW[7]}: nếu bật tiêm lỗi, luôn XOR mặt nạ \texttt{8'h55} tại vị trí 0 và \texttt{8'hA6} tại vị trí 14; nếu \texttt{SW[7]=1}, tiêm thêm lỗi thứ 3 với mặt nạ \texttt{8'h3C} tại vị trí 7.
    \item \textbf{\texttt{result\_checker.v} và \texttt{cycle\_counter.v}:} \texttt{u\_checker} chạy ở chế độ \texttt{MEM\_CHK} so sánh 11 byte giữa \texttt{u\_buf\_input} và \texttt{u\_buf\_output}, còn \texttt{u\_counter} đếm số chu kỳ 32-bit ở chế độ \textit{Pulse Mode} từ \texttt{S\_START} tới \texttt{decoder\_done}. Do các tín hiệu xuất của hai khối gỡ lỗi này chưa đấu ra LED ngoài, Quartus tối ưu bỏ phần logic dư lúc Place \& Route, nhưng vẫn được kiểm chứng đầy đủ trong \texttt{tb\_top}.
    \item \textbf{\texttt{reset\_sync.v}:} Mạch đồng bộ 2 tầng FF có thuộc tính \texttt{SYNCHRONIZER\_IDENTIFICATION FORCED\_IF\_ASYNCHRONOUS}, nhận \texttt{KEY[0]} để phát \texttt{rst\_n} hạ thấp bất đồng bộ tức thì và nhả đồng bộ theo xung nhịp 50\,MHz.
    \item \textbf{\texttt{input\_control.v}:} Đồng bộ 2 tầng FF cho \texttt{KEY[1]} và \texttt{SW[9:0]}, sau đó dùng bộ đếm 19 bit (\texttt{\$clog2(500\_000)}) lọc dội 10\,ms trước khi xuất xung 1 chu kỳ \texttt{confirm\_pulse} và chốt \texttt{sw\_latched}.
    \item \textbf{\texttt{display\_control.v} và \texttt{hex\_decoder.v}:} \texttt{display\_control} nối \texttt{core\_busy} ra \texttt{LEDR[0]} và chốt trạng thái \texttt{done\_latched}/\texttt{fail\_latched} ra \texttt{LEDR[1:2]} theo \texttt{clear\_status}; hai bộ \texttt{hex\_decoder} xuất mã LED 7 đoạn tích cực thấp với bit \texttt{DP = 1'b1}.
\end{itemize}

Đoạn mã Verilog dưới đây minh họa cơ chế đồng bộ 2 tầng và bộ lọc dội 10\,ms trong \texttt{input\_control.v}:

\begin{lstlisting}[language=Verilog]
always @(posedge clk or negedge rst_n) begin
    if (!rst_n) begin
        key_meta <= 1'b1; key_sync <= 1'b1;
        sw_meta  <= 10'b0; sw_sync <= 10'b0;
    end else begin
        key_meta <= key_confirm_n; key_sync <= key_meta;
        sw_meta  <= sw_raw;        sw_sync  <= sw_meta;
    end
end

always @(posedge clk or negedge rst_n) begin
    if (!rst_n) begin
        key_stable <= 1'b1; debounce_count <= {COUNTER_WIDTH{1'b0}};
        confirm_pulse <= 1'b0; sw_latched <= 10'b0;
    end else begin
        confirm_pulse <= 1'b0;
        if (key_sync == key_stable)
            debounce_count <= {COUNTER_WIDTH{1'b0}};
        else if (debounce_count == DEBOUNCE_CYCLES - 1) begin
            debounce_count <= {COUNTER_WIDTH{1'b0}};
            key_stable <= key_sync;
            if (key_sync == 1'b0) begin
                confirm_pulse <= 1'b1;
                sw_latched    <= sw_sync;
            end
        end else
            debounce_count <= debounce_count + 1'b1;
    end
end
\end{lstlisting}

\section{Kiểm thử mô phỏng (Testbench), Tổng hợp Quartus \& Thực nghiệm phần cứng}
\label{sec:verif-synth}

\subsection*{Kiểm thử mô phỏng tự động qua 10 Testbench}
Toàn bộ thiết kế RTL được kiểm chứng chéo với mô hình toán học chuẩn Python (\texttt{model/rs\_reference.py}, \texttt{model/gf256\_model.py}) trên ModelSim/QuestaSim thông qua tệp danh sách \texttt{sim/rtl\_files.f} và macro \texttt{sim/run.do} (độ phân giải \texttt{-t 1ps}). Mười testbench trong thư mục \texttt{tb/} (Bảng~\ref{tab:tb-summary}) đạt tổng cộng \textbf{105.175 lượt kiểm tra với 0 lỗi} (ghi chú: trong bản tóm tắt \texttt{README.md} có lỗi cộng nhầm dòng tổng thành 102.175).

\begin{table}[!htb]
    \centering
    \caption{Tổng hợp kết quả kiểm thử toàn bộ 10 Testbench RTL và hệ thống}
    \label{tab:tb-summary}
    \footnotesize
    \renewcommand{\arraystretch}{1.08}
    \begin{tabular}{l r c p{6.6cm}}
        \toprule
        \textbf{Tên Testbench} & \textbf{Số lượt} & \textbf{Kết quả} & \textbf{Phạm vi bao phủ kiểm thử} \\
        \midrule
        \texttt{tb\_gf256\_mul} & 72.820 & PASS (0 lỗi) & Đủ $256\times 256=65.536$ cặp, tính giao hoán/phân phối/kết hợp (2.000 mẫu) \\
        \texttt{tb\_gf256\_inv} & 1.538 & PASS (0 lỗi) & Vét cạn $0\dots 255$, $a\cdot a^{-1}=1$, song ánh, cờ \texttt{zero\_err/valid} \\
        \texttt{tb\_rs\_encoder} & 921 & PASS (0 lỗi) & Vector \texttt{00..0A}, 150 khối kèm backpressure, reset giữa chừng \\
        \texttt{tb\_syndrome} & 3.341 & PASS (0 lỗi) & 0 lỗi, $15\times 15$ mẫu 1 lỗi, 105 cặp 2 lỗi (gồm $S_0=0$), 160 mẫu $>T$ \\
        \texttt{tb\_berlekamp\_massey} & 10.437 & PASS (0 lỗi) & Vét cạn 1 lỗi \& 105 cặp $\times 4$ mẫu 2 lỗi, 150 mẫu $>T$, 200 ngẫu nhiên \\
        \texttt{tb\_error\_evaluator} & 2.971 & PASS (0 lỗi) & Đủ mẫu 1 lỗi, 105 cặp 2 lỗi ($\deg\Omega<L$), 300 mẫu ngẫu nhiên \\
        \texttt{tb\_chien\_search} & 2.834 & PASS (0 lỗi) & 15 mẫu 1 lỗi, 105 cặp 2 lỗi, nghiệm ngoài miền rút gọn ($k\ge 15\to\texttt{fail=1}$) \\
        \texttt{tb\_forney} & 5.525 & PASS (0 lỗi) & $15\times 37$ mẫu 1 lỗi, $105\times 3$ cặp 2 lỗi, bắt ngoại lệ $\Lambda'(z_i)=0$ \\
        \texttt{tb\_rs\_decoder} & 4.465 & PASS (0 lỗi) & Fast-path, đủ 1--2 lỗi, 300 mẫu $>T$ (Re-verify), 60 mẫu backpressure \\
        \texttt{tb\_top} & 323 & PASS (0 lỗi) & Debounce thực 500.000 chu kỳ ở 50\,MHz, lọc nhiễu 5 chu kỳ, 0/2/3 lỗi \\
        \midrule
        \textbf{Tổng cộng} & \textbf{105.175} & \textbf{100\% PASS} & \textbf{0 lỗi trên toàn bộ 10 cấp độ kiểm chứng} \\
        \bottomrule
    \end{tabular}
\end{table}

\subsection*{Kết quả tổng hợp Quartus Prime 25.1std.0 và phân tích thời gian tĩnh}
Quá trình biên dịch dự án \texttt{quartus/ReedSolomon.qpf} hoàn tất trong 48 giây và tạo tệp nạp \texttt{quartus/output\_files/ReedSolomon.sof}. Số liệu thực tế từ \texttt{ReedSolomon.fit.summary} và \texttt{ReedSolomon.fit.rpt} (Bảng~\ref{tab:quartus-res}) cho thấy thiết kế tiêu thụ \textbf{4.599 / 49.760 LEs (9\%)}, \textbf{2.035 thanh ghi (4\%)}, \textbf{346 / 3.110 LABs (11\%)}, \textbf{32 chân I/O (9\%)}, \textbf{0 M9K} và \textbf{0 bộ nhân $9\times 9$}. Con số thực tế này cao hơn ước tính sơ bộ ban đầu trong \texttt{Step\_to\_step.md} ($3\dots 6\%$ LEs, $800\dots 1.200$ FF) do toàn bộ bộ đệm 120-bit trong \texttt{rs\_decoder} giữ trên FF để đạt 0 trễ truy cập và cả 6 FSM dùng mã hóa \textit{One-Hot}.

\begin{table}[!htb]
    \centering
    \caption{Báo cáo tài nguyên phần cứng thực tế sau Place \& Route (Quartus Prime 25.1std.0)}
    \label{tab:quartus-res}
    \footnotesize
    \setlength{\tabcolsep}{4pt}
    \renewcommand{\arraystretch}{1.08}
    \begin{tabular}{l r r p{5.3cm}}
        \toprule
        \textbf{Module / Hạng mục} & \textbf{LEs} & \textbf{FFs} & \textbf{Tỷ lệ \& ghi chú kiến trúc} \\
        \midrule
        \textbf{Toàn hệ thống (\texttt{top\_de10\_lite})} & \textbf{4.599} & \textbf{2.035} & \textbf{9\% LEs}, 3.669 hàm tổ hợp (7\%), 346 LABs (11\%) \\
        \quad \texttt{u\_decoder} (\texttt{rs\_decoder}) & 4.026 & 1.658 & Chiếm 87,5\% tổng LEs (gồm bộ đệm 120-bit \& FSM) \\
        \qquad -- \texttt{u\_bm} (\texttt{berlekamp\_massey}) & 1.114 & 188 & Gồm \texttt{u\_binv} (208 LCs) và 10 khối \texttt{gf256\_mul} \\
        \qquad -- \texttt{u\_forney} (\texttt{forney}) & 831 & 207 & Gồm \texttt{u\_inv} (211 LCs) và mạch đánh giá Horner \\
        \qquad -- \texttt{u\_ec} / \texttt{u\_chien} / \texttt{u\_ee} & 434 / 391 / 382 & 366 / 78 / 104 & Sửa lỗi, Chien Search và Error Evaluator \\
        \qquad -- \texttt{u\_syndrome} (\texttt{syndrome}) & 85 & 71 & Tái sử dụng chung cho pha nhận và Re-verify \\
        \quad \texttt{u\_controller} / \texttt{u\_encoder} & 186 / 101 & 26 / 53 & Điều phối trung tâm và LFSR mã hóa \\
        \quad \texttt{u\_buf\_output} / \texttt{u\_buf\_input} & 147 / 137 & 124 / 101 & Bộ đệm thanh ghi 15 byte và 11 byte \\
        \quad Ngoại vi (\texttt{input, injector, display, reset}) & 106 & 73 & Đồng bộ hóa, chống dội, tiêm lỗi và hiển thị \\
        \midrule
        \textbf{I/O / RAM M9K / DSP 9$\times$9} & \textbf{32 / 360} & \textbf{0 / 0} & 13 Input, 19 Output (9\%); 0 M9K và 0 DSP \\
        \bottomrule
    \end{tabular}
\end{table}

Phân tích thời gian tĩnh TimeQuest (\texttt{ReedSolomon.sta.rpt}, Bảng~\ref{tab:sta-results}) xác nhận mạch đóng timing an toàn ở 50\,MHz tại mọi góc nhiệt độ/điện áp, đạt $F_{\max} = 51{,}02\text{ MHz}$ và \textit{Setup Slack} $= +0{,}398\text{ ns}$ (\texttt{TNS = 0,000 ns}) ở góc Slow $85^{\circ}\text{C}$, cùng độ tin cậy siêu ổn định $\text{MTBF} \ge 10^9$ năm trên 12 chuỗi đồng bộ. Đường găng thời gian nằm trên chuỗi tổ hợp \texttt{gf256\_inv} nối tiếp \texttt{gf256\_mul} trong \texttt{berlekamp\_massey} và \texttt{forney}.

\begin{table}[!htb]
    \centering
    \caption{Kết quả phân tích thời gian tĩnh (TimeQuest STA) với xung nhịp 50\,MHz}
    \label{tab:sta-results}
    \footnotesize
    \setlength{\tabcolsep}{4pt}
    \renewcommand{\arraystretch}{1.1}
    \begin{tabular}{l c c c c c}
        \toprule
        \textbf{Góc phân tích (Corner)} & \textbf{$F_{\max}$ (MHz)} & \textbf{Setup Slack} & \textbf{Hold Slack} & \textbf{Rec. / Rem. (ns)} & \textbf{Min Pulse (ns)} \\
        \midrule
        Slow 1200\,mV 85\,$^{\circ}\text{C}$ & 51,02 & $+0{,}398\text{ ns}$ & $+0{,}341\text{ ns}$ & $+16{,}638$ / $+2{,}350$ & $+9{,}624$ \\
        Slow 1200\,mV 0\,$^{\circ}\text{C}$ & 55,42 & $+1{,}955\text{ ns}$ & $+0{,}305\text{ ns}$ & $+17{,}030$ / $+2{,}093$ & $+9{,}655$ \\
        Fast 1200\,mV 0\,$^{\circ}\text{C}$ & -- & $+11{,}687\text{ ns}$ & $+0{,}148\text{ ns}$ & $+18{,}541$ / $+1{,}109$ & $+9{,}741$ \\
        \bottomrule
    \end{tabular}
\end{table}

\subsection*{Kịch bản thực nghiệm kiểm chứng trên phần cứng DE10-Lite}
Sau khi nạp \texttt{ReedSolomon.sof} qua USB-Blaster, 4 kịch bản thực nghiệm với chuỗi mẫu \texttt{0x00..0x0A} khẳng định đúng hoạt động phần cứng:
(1) \textbf{0 lỗi (\texttt{SW[9]=0}):} qua \textit{Fast-path}, \texttt{LEDR[2:0] = 3'b010}, kiểm tra địa chỉ $11\dots 14$ qua \texttt{SW[3:0]} hiển thị đúng 4 byte parity \texttt{C0, 8F, 28, 6C};
(2) \textbf{2 lỗi (\texttt{SW[9]=1, SW[7]=0}):} tiêm \texttt{0x55} tại vị trí 0 và \texttt{0xA6} tại vị trí 14, bộ giải mã tự sửa sạch và xuất \texttt{LEDR = 3'b010}, kiểm tra byte 0 và 14 phục hồi đúng \texttt{00} và \texttt{6C};
(3) \textbf{3 lỗi vượt ngưỡng (\texttt{SW[9]=1, SW[7]=1}):} tiêm thêm \texttt{0x3C} tại vị trí 7, mạch bật \texttt{LEDR = 3'b110} (báo Fail) và byte 0 hiển thị \texttt{55} ($0\text{x00}\oplus 0\text{x55}$), chứng minh giữ nguyên dữ liệu nhận thay vì xuất kết quả sửa sai;
(4) \textbf{Reset (\texttt{KEY[0]}):} nhấn ở thời điểm bất kỳ đưa ngay \texttt{LEDR} về \texttt{3'b000}.


% ============================================================================
% PHAN CUOI: KET LUAN VA TAI LIEU THAM KHAO
% ============================================================================
\chapter*{Kết luận và Hướng phát triển}
\phantomsection
\addcontentsline{toc}{chapter}{Kết luận và Hướng phát triển}
\markboth{Kết luận và Hướng phát triển}{}
\label{ch:conclusion}

Đồ án đã hoàn thành thiết kế, kiểm chứng mô phỏng và triển khai thực tế hệ thống mã hóa -- giải mã Reed--Solomon $\mathrm{RS}(15,11)$ trên kit FPGA Terasic DE10-Lite (\texttt{10M50DAF484C7G}). Ba kết quả trọng tâm đã đạt được gồm: (1) kiến trúc RTL 21 tệp Verilog phân tầng chặt chẽ theo giao thức \texttt{valid}/\texttt{ready} có backpressure; (2) xác minh đúng đắn $100\%$ qua 10 testbench với tổng \textbf{105.175 lượt kiểm tra (0 lỗi)} đối chiếu mô hình Python; và (3) tổng hợp thành công trên Quartus Prime \texttt{25.1std.0} chỉ tiêu tốn \textbf{4.599 LEs (9\%)}, \textbf{2.035 FF (4\%)}, $0\%$ M9K/DSP và đóng timing ở $F_{\max} = 51{,}02\text{ MHz}$ tại tần số bo mạch 50\,MHz.

\subsection*{Đánh giá ưu điểm và hạn chế}
Về ưu điểm, hệ thống tích hợp cơ chế \textit{Fast-path} (giảm độ trễ sau byte nhận cuối xuống 2 chu kỳ khi $S=0$), cơ chế kiểm tra hai lớp kết hợp tái sử dụng duy nhất một khối \texttt{u\_syndrome} cho pha Re-verify nhằm loại bỏ hoàn toàn giải mã sai thầm lặng khi số lỗi $>2$, cùng tầng đồng bộ hóa ngoại vi đạt $\text{MTBF} \ge 10^9$ năm. Tuy vậy, thiết kế hiện tại còn ba hạn chế chính: chưa hỗ trợ giải mã xóa (\textit{Erasure}); xử lý nối tiếp 1 symbol/chu kỳ và đường găng tổ hợp \texttt{gf256\_inv} + \texttt{gf256\_mul} giữ $F_{\max}$ ở mức $51{,}02\text{ MHz}$ (biên setup $+0{,}398\text{ ns}$); tham số $(N,K,R)$ cấu hình tĩnh và hai khối giám sát \texttt{result\_checker}, \texttt{cycle\_counter} chưa nối ra chân giao tiếp ngoài nên bị tối ưu lược bỏ lúc Place \& Route.

\subsection*{Hướng phát triển tiếp theo}
Đề tài có thể mở rộng trực tiếp theo ba hướng kỹ thuật cụ thể:
\begin{enumerate}
    \item \textbf{Pipeline hóa khối BM/Forney để nâng tần số trên 100\,MHz:} Chèn thanh ghi đệm sau LUT nghịch đảo hoặc áp dụng thuật toán Berlekamp--Massey không phép nghịch đảo (\textit{Inversionless BM}) nhằm cắt đường găng tổ hợp và tăng thông lượng xử lý khung liên tiếp.
    \item \textbf{Hỗ trợ giải mã Xóa (\textit{Errors-and-Erasures}) và xuất telemetry qua UART:} Nâng giới hạn sửa lỗi lên $2t + \rho \le 4$ và sử dụng hai chân \texttt{PIN\_V10}/\texttt{PIN\_W10} trên header \texttt{JP1} để truyền dữ liệu từ mã, số chu kỳ đếm được từ \texttt{cycle\_counter} và kết quả \texttt{result\_checker} lên máy tính.
    \item \textbf{Mở rộng sang chuẩn đầy đủ $\mathrm{RS}(255,239)$ dùng bộ nhớ M9K:} Chuyển bộ đệm thanh ghi sang RAM nhúng \textbf{M9K} sẵn có trên MAX 10 và song song hóa các nhánh đánh giá Chien Search để đáp ứng các chuẩn truyền thông tốc độ cao.
\end{enumerate}

\begin{thebibliography}{99}

\bibitem{reed1960}
I.~S.~Reed và G.~Solomon, ``Polynomial Codes over Certain Finite Fields,'' \textit{Journal of the Society for Industrial and Applied Mathematics (SIAM)}, tập~8, số~2, tr.~300--304, tháng~6/1960.
\addcontentsline{toc}{chapter}{Tài liệu tham khảo}

\bibitem{berlekamp1968}
E.~R.~Berlekamp, \textit{Algebraic Coding Theory}, McGraw-Hill Book Company, New York, 1968.

\bibitem{massey1969}
J.~L.~Massey, ``Shift-Register Synthesis and BCH Decoding,'' \textit{IEEE Transactions on Information Theory}, tập~IT-15, số~1, tr.~122--127, tháng~1/1969.

\bibitem{chien1964}
R.~T.~Chien, ``Cyclic Decoding Procedures for Bose--Chaudhuri--Hocquenghem Codes,'' \textit{IEEE Transactions on Information Theory}, tập~IT-10, số~4, tr.~357--363, tháng~10/1964.

\bibitem{forney1965}
G.~D.~Forney, ``On Decoding BCH Codes,'' \textit{IEEE Transactions on Information Theory}, tập~IT-11, số~4, tr.~549--557, tháng~10/1965.

\bibitem{lincostello2004}
S.~Lin và D.~J.~Costello, \textit{Error Control Coding: Fundamentals and Applications}, ấn bản lần~2, Pearson Prentice Hall, Upper Saddle River, NJ, 2004.

\bibitem{morelos2006}
R.~H.~Morelos-Zaragoza, \textit{The Art of Error Correcting Coding}, ấn bản lần~2, John Wiley \& Sons, Chichester, UK, 2006.

\bibitem{intelmax10}
Intel Corporation, \textit{Intel MAX 10 FPGA Device Datasheet}, Tài liệu kỹ thuật \texttt{M10-DATASHEET}, Intel Programmable Solutions Group, 2022.

\bibitem{de10lite}
Terasic Inc., \textit{DE10-Lite User Manual}, Tài liệu hướng dẫn sử dụng kit Terasic DE10-Lite (\texttt{10M50DAF484C7G}), Hsinchu, Đài Loan, 2020.

\bibitem{quartussta}
Intel Corporation, \textit{Intel Quartus Prime Standard Edition User Guide: Timing Analyzer}, Phiên bản \texttt{25.1std.0}, Intel FPGA Documentation, 2025.

\bibitem{mathworksrs}
The MathWorks, Inc., \textit{HDL Reed--Solomon Encoder and Decoder Reference Architecture}, \textit{Communications Toolbox HDL Support Documentation}, Natick, MA, 2024.

\end{thebibliography}

\end{document}
